% Generated draft: requires mathematical and source review.
\documentclass[11pt]{article}
\usepackage[margin=25mm]{geometry}
\usepackage{fontspec}
\setmainfont{Latin Modern Roman}
\usepackage{amsmath,amssymb,mathtools,unicode-math}
\setmathfont{Latin Modern Math}
\usepackage{enumitem,needspace,xurl,hyperref}
\setlength{\parindent}{0pt}
\setlength{\parskip}{5pt}
\setlength{\emergencystretch}{3em}
\newcommand{\N}{\mathbb{N}}
\newcommand{\Z}{\mathbb{Z}}
\newcommand{\Q}{\mathbb{Q}}
\newcommand{\R}{\mathbb{R}}
\newcommand{\C}{\mathbb{C}}
\newcommand{\F}{\mathbb{F}}
\newcommand{\A}{\mathbb{A}}
\newcommand{\PP}{\mathbb{P}}
\newcommand{\E}{\mathbb{E}}
\newcommand{\eps}{\varepsilon}
\newcommand{\dd}{\,\mathrm{d}}
\begin{document}
\section*{UnsolvedMath ID: 30000502}
\textbf{Source problem code:} catalogue entry
\subsection{Definitions and mathematical statement}
The source asks what information about ideals in an Iwasawa algebra can be obtained from its Lie algebra of inner derivations and the induced Poisson structure on its associated graded algebra. It specifies neither the group/coefficient setting nor a particular ideal-theoretic conclusion, so this is a research direction rather than a single quantified classification theorem.

For an associative algebra \(A\), an inner derivation is \(\operatorname{ad}_a(x)=[a,x]=ax-xa\). The identities
\[
 [\operatorname{ad}_a,\operatorname{ad}_b]=\operatorname{ad}_{[a,b]}
\]
give the inner derivations a Lie-algebra structure, with central elements acting trivially. Every two-sided ideal \(I\) satisfies \(\operatorname{ad}_a(I)\subseteq I\) for all \(a\). This supplies the basic connection to ideals but not a classification.

To describe the graded mechanism conditionally, suppose an increasing multiplicative filtration \(F_\bullet A\) has been specified, and assume for some fixed \(\delta>0\) that
\[
 [F_iA,F_jA]\subseteq F_{i+j-\delta}A
\]
for all relevant indices. Then the corresponding leading commutator gives a bracket of degree \(-\delta\) on \(\operatorname{gr}_F A\), by
\[
 \{\sigma_i(a),\sigma_j(b)\}
 =\sigma_{i+j-\delta}([a,b]),
\]
with its induced skew, Jacobi and Leibniz identities. This is one conventional filtered-algebra setup, not an assertion that a particular Iwasawa filtration in the source satisfies it. Other filtrations and normalizations change the degree and construction. An induced filtered two-sided ideal has a graded ideal stable under the corresponding leading commutators, when this setup applies.

The requested information could concern constraints on graded ideals, prime or other ideals, their supports, or lifting graded data back to the completed algebra, but the source chooses none of these precise targets. In particular, the existence of a Poisson-stable graded ideal does not by itself assert a unique lift, a classification of ideals or preservation of primeness under passage to the associated graded ring. The intended ideal class, closure conditions, filtration and level of conclusions must be supplied before making any such claim.

The literature triage checked 2026-08-21 describes the extraction as a transition question with no recoverable object or conclusion and labels it open rather than discarding it. That description does not match the explicit Iwasawa/inner-derivation/graded-Poisson objects in both current statement fields, though the desired conclusion remains broad. The exact-statement verification and open label are retained as source metadata, not independently verified status.
\subsection{Short English statement}
Investigate how inner derivations and their leading Poisson brackets constrain ideals of the intended Iwasawa algebra. The filtration, coefficient/group hypotheses and desired ideal-theoretic result are missing. A graded Poisson description is not automatically an ideal classification or a lifting theorem.
\subsection{Sources}
\begin{itemize}
\item Local UnsolvedMath registry, permanent ID 30000502, source code catalogue entry, ``Ideals in Iwasawa Algebras from Inner Derivations,'' including the unchanged original question, exact-statement verification and conflicting extraction assessment.
\item \emph{Pro-\(p\) Extensions of Global Fields and pro-\(p\) Groups (2007)}, Oberwolfach Reports, pp.\ 1463--1536, DOI 10.4171/owr/2006/25: \url{https://doi.org/10.4171/owr/2006/25}. This is the sole supplied literature reference, labelled ``original OWR report.'' The citation uses 2007, while the DOI and proposed-year metadata use 2006.
\item Literature triage checked 2026-08-21; local record updated 2026-09-16. No external source has been visited or independently checked here.
\end{itemize}
\end{document}
